\IfFileExists{stacks-project.cls}{%
\documentclass{stacks-project}
}{%
\documentclass{amsart}
}

% For dealing with references we use the comment environment
\usepackage{verbatim}
\newenvironment{reference}{\comment}{\endcomment}
%\newenvironment{reference}{}{}
\newenvironment{slogan}{\comment}{\endcomment}
\newenvironment{history}{\comment}{\endcomment}

% For commutative diagrams we use Xy-pic
\usepackage[all]{xy}

% We use 2cell for 2-commutative diagrams.
\xyoption{2cell}
\UseAllTwocells

% We use multicol for the list of chapters between chapters
\usepackage{multicol}

% This is generally recommended for better output
\usepackage{lmodern}
\usepackage[T1]{fontenc}

% For cross-file-references
\usepackage{xr-hyper}

% Package for hypertext links:
\usepackage{hyperref}

% For any local file, say "hello.tex" you want to link to please
% use \externaldocument[hello-]{hello}
\externaldocument[introduction-]{introduction}
\externaldocument[conventions-]{conventions}
\externaldocument[sets-]{sets}
\externaldocument[categories-]{categories}
\externaldocument[topology-]{topology}
\externaldocument[sheaves-]{sheaves}
\externaldocument[sites-]{sites}
\externaldocument[stacks-]{stacks}
\externaldocument[fields-]{fields}
\externaldocument[algebra-]{algebra}
\externaldocument[brauer-]{brauer}
\externaldocument[homology-]{homology}
\externaldocument[derived-]{derived}
\externaldocument[simplicial-]{simplicial}
\externaldocument[more-algebra-]{more-algebra}
\externaldocument[smoothing-]{smoothing}
\externaldocument[modules-]{modules}
\externaldocument[sites-modules-]{sites-modules}
\externaldocument[injectives-]{injectives}
\externaldocument[cohomology-]{cohomology}
\externaldocument[sites-cohomology-]{sites-cohomology}
\externaldocument[dga-]{dga}
\externaldocument[dpa-]{dpa}
\externaldocument[sdga-]{sdga}
\externaldocument[hypercovering-]{hypercovering}
\externaldocument[schemes-]{schemes}
\externaldocument[constructions-]{constructions}
\externaldocument[properties-]{properties}
\externaldocument[morphisms-]{morphisms}
\externaldocument[coherent-]{coherent}
\externaldocument[divisors-]{divisors}
\externaldocument[limits-]{limits}
\externaldocument[varieties-]{varieties}
\externaldocument[topologies-]{topologies}
\externaldocument[descent-]{descent}
\externaldocument[perfect-]{perfect}
\externaldocument[more-morphisms-]{more-morphisms}
\externaldocument[flat-]{flat}
\externaldocument[groupoids-]{groupoids}
\externaldocument[more-groupoids-]{more-groupoids}
\externaldocument[etale-]{etale}
\externaldocument[chow-]{chow}
\externaldocument[intersection-]{intersection}
\externaldocument[pic-]{pic}
\externaldocument[weil-]{weil}
\externaldocument[adequate-]{adequate}
\externaldocument[dualizing-]{dualizing}
\externaldocument[duality-]{duality}
\externaldocument[discriminant-]{discriminant}
\externaldocument[derham-]{derham}
\externaldocument[local-cohomology-]{local-cohomology}
\externaldocument[algebraization-]{algebraization}
\externaldocument[curves-]{curves}
\externaldocument[resolve-]{resolve}
\externaldocument[models-]{models}
\externaldocument[functors-]{functors}
\externaldocument[equiv-]{equiv}
\externaldocument[pione-]{pione}
\externaldocument[etale-cohomology-]{etale-cohomology}
\externaldocument[proetale-]{proetale}
\externaldocument[relative-cycles-]{relative-cycles}
\externaldocument[more-etale-]{more-etale}
\externaldocument[trace-]{trace}
\externaldocument[crystalline-]{crystalline}
\externaldocument[spaces-]{spaces}
\externaldocument[spaces-properties-]{spaces-properties}
\externaldocument[spaces-morphisms-]{spaces-morphisms}
\externaldocument[decent-spaces-]{decent-spaces}
\externaldocument[spaces-cohomology-]{spaces-cohomology}
\externaldocument[spaces-limits-]{spaces-limits}
\externaldocument[spaces-divisors-]{spaces-divisors}
\externaldocument[spaces-over-fields-]{spaces-over-fields}
\externaldocument[spaces-topologies-]{spaces-topologies}
\externaldocument[spaces-descent-]{spaces-descent}
\externaldocument[spaces-perfect-]{spaces-perfect}
\externaldocument[spaces-more-morphisms-]{spaces-more-morphisms}
\externaldocument[spaces-flat-]{spaces-flat}
\externaldocument[spaces-groupoids-]{spaces-groupoids}
\externaldocument[spaces-more-groupoids-]{spaces-more-groupoids}
\externaldocument[bootstrap-]{bootstrap}
\externaldocument[spaces-pushouts-]{spaces-pushouts}
\externaldocument[spaces-chow-]{spaces-chow}
\externaldocument[groupoids-quotients-]{groupoids-quotients}
\externaldocument[spaces-more-cohomology-]{spaces-more-cohomology}
\externaldocument[spaces-simplicial-]{spaces-simplicial}
\externaldocument[spaces-duality-]{spaces-duality}
\externaldocument[formal-spaces-]{formal-spaces}
\externaldocument[restricted-]{restricted}
\externaldocument[spaces-resolve-]{spaces-resolve}
\externaldocument[formal-defos-]{formal-defos}
\externaldocument[defos-]{defos}
\externaldocument[cotangent-]{cotangent}
\externaldocument[examples-defos-]{examples-defos}
\externaldocument[algebraic-]{algebraic}
\externaldocument[examples-stacks-]{examples-stacks}
\externaldocument[stacks-sheaves-]{stacks-sheaves}
\externaldocument[criteria-]{criteria}
\externaldocument[artin-]{artin}
\externaldocument[quot-]{quot}
\externaldocument[stacks-properties-]{stacks-properties}
\externaldocument[stacks-morphisms-]{stacks-morphisms}
\externaldocument[stacks-limits-]{stacks-limits}
\externaldocument[stacks-cohomology-]{stacks-cohomology}
\externaldocument[stacks-perfect-]{stacks-perfect}
\externaldocument[stacks-introduction-]{stacks-introduction}
\externaldocument[stacks-more-morphisms-]{stacks-more-morphisms}
\externaldocument[stacks-geometry-]{stacks-geometry}
\externaldocument[moduli-]{moduli}
\externaldocument[moduli-curves-]{moduli-curves}
\externaldocument[examples-]{examples}
\externaldocument[exercises-]{exercises}
\externaldocument[guide-]{guide}
\externaldocument[desirables-]{desirables}
\externaldocument[coding-]{coding}
\externaldocument[obsolete-]{obsolete}
\externaldocument[fdl-]{fdl}
\externaldocument[index-]{index}

% Theorem environments.
%
\theoremstyle{plain}
\newtheorem{theorem}[subsection]{Theorem}
\newtheorem{proposition}[subsection]{Proposition}
\newtheorem{lemma}[subsection]{Lemma}

\theoremstyle{definition}
\newtheorem{definition}[subsection]{Definition}
\newtheorem{example}[subsection]{Example}
\newtheorem{exercise}[subsection]{Exercise}
\newtheorem{situation}[subsection]{Situation}

\theoremstyle{remark}
\newtheorem{remark}[subsection]{Remark}
\newtheorem{remarks}[subsection]{Remarks}

\numberwithin{equation}{subsection}

% Macros
%
\def\lim{\mathop{\mathrm{lim}}\nolimits}
\def\colim{\mathop{\mathrm{colim}}\nolimits}
\def\Spec{\mathop{\mathrm{Spec}}}
\def\Hom{\mathop{\mathrm{Hom}}\nolimits}
\def\Ext{\mathop{\mathrm{Ext}}\nolimits}
\def\SheafHom{\mathop{\mathcal{H}\!\mathit{om}}\nolimits}
\def\SheafExt{\mathop{\mathcal{E}\!\mathit{xt}}\nolimits}
\def\Sch{\mathit{Sch}}
\def\Mor{\mathop{\mathrm{Mor}}\nolimits}
\def\Ob{\mathop{\mathrm{Ob}}\nolimits}
\def\Sh{\mathop{\mathit{Sh}}\nolimits}
\def\NL{\mathop{N\!L}\nolimits}
\def\CH{\mathop{\mathrm{CH}}\nolimits}
\def\proetale{{pro\text{-}\acute{e}tale}}
\def\etale{{\acute{e}tale}}
\def\QCoh{\mathit{QCoh}}
\def\Ker{\mathop{\mathrm{Ker}}}
\def\Im{\mathop{\mathrm{Im}}}
\def\Coker{\mathop{\mathrm{Coker}}}
\def\Coim{\mathop{\mathrm{Coim}}}

% Boxtimes
%
\DeclareMathSymbol{\boxtimes}{\mathbin}{AMSa}{"02}

%
% Macros for moduli stacks/spaces
%
\def\QCohstack{\mathcal{QC}\!\mathit{oh}}
\def\Cohstack{\mathcal{C}\!\mathit{oh}}
\def\Spacesstack{\mathcal{S}\!\mathit{paces}}
\def\Quotfunctor{\mathrm{Quot}}
\def\Hilbfunctor{\mathrm{Hilb}}
\def\Curvesstack{\mathcal{C}\!\mathit{urves}}
\def\Polarizedstack{\mathcal{P}\!\mathit{olarized}}
\def\Complexesstack{\mathcal{C}\!\mathit{omplexes}}
% \Pic is the operator that assigns to X its picard group, usage \Pic(X)
% \Picardstack_{X/B} denotes the Picard stack of X over B
% \Picardfunctor_{X/B} denotes the Picard functor of X over B
\def\Pic{\mathop{\mathrm{Pic}}\nolimits}
\def\Picardstack{\mathcal{P}\!\mathit{ic}}
\def\Picardfunctor{\mathrm{Pic}}
\def\Deformationcategory{\mathcal{D}\!\mathit{ef}}
\usepackage{cleveref}
\crefname{originallemma}{original lemma}{original lemmas}
\Crefname{originallemma}{Original lemma}{Original lemmas}
\makeatletter
\newlabel{lemma-finite-flat-over-almost-dense-subscheme}{{0ABA}{source}{}{sourcepointer-1}{}}
\newlabel{lemma-finite-flat-over-almost-dense-subscheme@cref}{{[originallemma][0][]0ABA}{[1][0][]source}{}{}{}}
\makeatother
\crefname{originallemma}{original lemma}{original lemmas}
\Crefname{originallemma}{Original lemma}{Original lemmas}
\makeatletter
\newlabel{coherent-lemma-image-affine-finite-morphism-affine-Noetherian}{{01YQ}{source}{}{sourcepointer-2}{}}
\newlabel{coherent-lemma-image-affine-finite-morphism-affine-Noetherian@cref}{{[originallemma][0][]01YQ}{[1][0][]source}{}{}{}}
\makeatother
\crefname{originallemma}{original lemma}{original lemmas}
\Crefname{originallemma}{Original lemma}{Original lemmas}
\makeatletter
\newlabel{properties-lemma-maximal-points-affine}{{01ZV}{source}{}{sourcepointer-3}{}}
\newlabel{properties-lemma-maximal-points-affine@cref}{{[originallemma][0][]01ZV}{[1][0][]source}{}{}{}}
\makeatother
\crefname{originallemma}{original lemma}{original lemmas}
\Crefname{originallemma}{Original lemma}{Original lemmas}
\makeatletter
\newlabel{properties-lemma-locally-Noetherian-quasi-separated}{{01OY}{source}{}{sourcepointer-4}{}}
\newlabel{properties-lemma-locally-Noetherian-quasi-separated@cref}{{[originallemma][0][]01OY}{[1][0][]source}{}{}{}}
\makeatother
\crefname{originallemma}{original lemma}{original lemmas}
\Crefname{originallemma}{Original lemma}{Original lemmas}
\makeatletter
\newlabel{groupoids-lemma-find-invariant-affine}{{03JE}{source}{}{sourcepointer-5}{}}
\newlabel{groupoids-lemma-find-invariant-affine@cref}{{[originallemma][0][]03JE}{[1][0][]source}{}{}{}}
\makeatother
\title{R-invariant affine thickening pairs around generic orbit points for locally Noetherian U with finite source and target}
\author{The Stacks Project Authors}
\date{stacks-project:more-groupoids.tex:local-snapshot-20261012; GNU FDL 1.2 or later}
\begin{document}
\maketitle
\setcounter{subsection}{1}
\noindent\textit{Editorial note.} All delivered mathematical text is the author's own, copied from the pinned Stacks Project chapter file more-groupoids.tex; no statement or proof step was rewritten, reordered or removed. The counted span is the author environment and its own proof environment. References inside the span are handed to the assembler as navigation stubs that carry the canonical tag of the cited tag only; whether each of them is load-bearing is recorded on the stub itself.\par\medskip
\noindent\textit{Allowed standard prerequisites (editorial).}\begin{itemize}
\item The R-invariant thickening pair W' inside W with finite locally free restriction and the covering property (more-groupoids chapter, Lemma lemma-finite-flat-over-almost-dense-subscheme, canonical tag 0ABA; collected on this lane as L22511). cited at proof lines 2861-2862 and used through lines 2863-2874. Provides the pair W' inside W whose properties drive the orbit-containment and the replacement reduction; its covering (2) puts some u\_\{ji\} in the closure of W.
\item Images of affine schemes by finite morphisms are affine, Noetherian setting (cohomology of schemes chapter, Lemma coherent-lemma-image-affine-finite-morphism-affine-Noetherian, canonical tag 01YQ). cited at proof lines 2868-2872. Reduces the task to finding an R-invariant affine open of W': the corresponding V inside W is then affine since W' to W is finite.
\item Affine opens containing finitely many maximal-type points of a quasi-separated scheme (properties chapter, Lemma properties-lemma-maximal-points-affine, canonical tag 01ZV). cited at proof lines 2877-2878. Finds an affine open containing the finitely many orbit points in the restricted, quasi-separated situation.
\item A locally Noetherian scheme is quasi-separated (properties chapter, Lemma properties-lemma-locally-Noetherian-quasi-separated, canonical tag 01OY). cited at proof lines 2879-2880. Justifies the applicability of the maximal-points lemma to the locally Noetherian U (and its opens).
\item Existence of R-invariant affine opens containing finitely many points for groupoids with finite locally free structure maps (groupoids chapter, Lemma groupoids-lemma-find-invariant-affine, canonical tag 03JE). cited at proof lines 2881-2883. Produces the final R-invariant affine open containing the orbit points in the finite locally free restriction.
\end{itemize}\medskip
\ifdefined\hypertarget\hypertarget{sourcepointer-1}{}\fi
\noindent\textit{Source reference.} Original lemma 0ABA: sources/stacks-chow/tags.txt, source lines 13013--13013. Cited at proof lines 2861-2862: supplies the R-invariant pair W' inside W, its covering property (2), and the finite locally free restriction used in the reduction.\par\medskip
\ifdefined\hypertarget\hypertarget{sourcepointer-2}{}\fi
\noindent\textit{Source reference.} Original lemma 01YQ: sources/stacks-chow/tags.txt, source lines 2565--2565. Cited at proof lines 2868-2872: reduces the construction to an R-invariant affine open of W', the thickening V being affine by finiteness.\par\medskip
\ifdefined\hypertarget\hypertarget{sourcepointer-3}{}\fi
\noindent\textit{Source reference.} Original lemma 01ZV: sources/stacks-chow/tags.txt, source lines 2610--2610. Cited at proof lines 2877-2878: finds an affine open containing the finitely many orbit points.\par\medskip
\ifdefined\hypertarget\hypertarget{sourcepointer-4}{}\fi
\noindent\textit{Source reference.} Original lemma 01OY: sources/stacks-chow/tags.txt, source lines 2211--2211. Cited at proof lines 2879-2880: a locally Noetherian scheme is quasi-separated, the hypothesis needed for the maximal-points lemma.\par\medskip
\ifdefined\hypertarget\hypertarget{sourcepointer-5}{}\fi
\noindent\textit{Source reference.} Original lemma 03JE: sources/stacks-chow/tags.txt, source lines 4636--4636. Cited at proof lines 2881-2883: produces the R-invariant affine open containing the orbit points for a groupoid with finite locally free structure maps.\par\medskip
\begin{lemma}
\label{lemma-invariant-affine-open-around-generic-point-Noetherian}
Let $(U, R, s, t, c)$ be a groupoid scheme over a scheme $S$.
Assume $s, t$ finite, $U$ is locally Noetherian, and $u_1, \ldots, u_m \in U$
points whose orbits consist of generic points of irreducible
components of $U$. Then there exist $R$-invariant subschemes
$V' \subset V \subset U$ such that
\begin{enumerate}
\item $u_1, \ldots, u_m \in V'$,
\item $V$ is open in $U$,
\item $V'$ and $V$ are affine,
\item $V' \subset V$ is a thickening,
\item the morphisms $s', t'$ of the restriction $(V', R', s', t', c')$
are finite locally free.
\end{enumerate}
\end{lemma}

\begin{proof}
Let $\{u_{j1}, \ldots, u_{jn_j}\}$ be the orbit of $u_j$.
Let $W' \subset W \subset U$ be as in
Lemma \ref{lemma-finite-flat-over-almost-dense-subscheme}.
Since $U = t(s^{-1}(\overline{W}))$ we see that at least
one $u_{ji} \in \overline{W}$. Since $u_{ji}$ is a generic point
of an irreducible component and $U$ locally Noetherian,
this implies that $u_{ji} \in W$. Since $W$ is $R$-invariant, we
conclude that $u_j \in W$ and in fact the whole orbit is contained in $W$.
By Cohomology of Schemes, Lemma
\ref{coherent-lemma-image-affine-finite-morphism-affine-Noetherian}
it suffices to find an $R$-invariant affine open subscheme $V'$
of $W'$ containing $u_1, \ldots, u_m$ (because then we can let $V \subset W$
be the corresponding open subscheme which will be affine).
Thus we may replace $(U, R, s, t, c)$
by the restriction $(W', R', s', t', c')$ to $W'$.
In other words, we may assume we have a groupoid scheme $(U, R, s, t, c)$
whose morphisms $s$ and $t$ are finite locally free.
By Properties, Lemma \ref{properties-lemma-maximal-points-affine}
we can find an affine open containing $\{u_{ij}\}$
(a locally Noetherian scheme is quasi-separated by
Properties, Lemma \ref{properties-lemma-locally-Noetherian-quasi-separated}).
Finally, we can apply
Groupoids, Lemma \ref{groupoids-lemma-find-invariant-affine}
to conclude.
\end{proof}

\end{document}
